\chapter{Уншсан судалгаа ба гаргасан дүгнэлт}
\label{ch:literature}

\section{Сэдэв, чиглэл сонгосон нь (10-р сар, 1-р долоо хоног)}
Анх "ханш өсөх магадлал 20\%, буурах магадлал 80\%" гэсэн нэг тоо гаргадаг загвар хийх санаа байсан. Гэвч ийм нэг тоо
арилжаачинд найдваргүй байдалыг төрүүлдэг. Хэрэвээ тэр хүн энэ үзүүлэлтэд итгэн дагаж арилжахад алдагдал хүлээвэл энэ нь олон асуудалыг үүсгэнэ. Иймд ханшийн боломжит хэд хэдэн замыг (хувилбарыг) тус бүрийн
магадлалтай нь харуулах чиглэлийг сонгосон. Оролтод үнийн бүтэц (лаа) болон эдийн засгийн мэдээг ашиглахаар тогтсон.

\section{FLF-LSTM судалгаа (10-р сар, 2-р долоо хоног)}
Ahmed нарын FLF-LSTM загвар \cite{ahmed2020flf} нь EUR/USD-ийн 4 цагийн лааны нээлт, дээд, доод, хаалтын үнээс
(OHLC) дараагийн лааг LSTM \cite{hochreiter1997lstm}-ээр таамагладаг. Судлаачид лааны дунд цэгийн зөрүүг алдагдлын
функцдээ нэмснээр энгийн LSTM-ээс 10.96\% бага алдаатай болсон гэж мэдээлсэн. Энэ нь лааны бүтэц шугаман графикаас
(зөвхөн хаалтын үнэ) илүү мэдээлэл агуулдаг гэсэн санааг дэмжиж байсан тул судалгааг дэлгэрэнгүй шалгасан.

Судлаачдын GitHub-т нийтэлсэн өгөгдөл, таамаглалын файлуудыг уншиж шалгахад дараах зүйл илэрсэн:
\begin{itemize}
    \item Тэдний яг ижил test өгөгдөл дээр "дараагийн лаа = одоогийн лаа" гэсэн хамгийн энгийн таамаглал дунджаар
          14.8 пип алдаатай. Харин нийтлэлд FLF-LSTM-ийн алдааг 21.9 пип гэж мэдээлсэн. Өөрөөр хэлбэл санал болгосон
          загвар энгийн таамаглалаас муу байна. Нийтлэлд ийм харьцуулалт хийгээгүй.
    \item Хадгалсан таамаглалуудын 73--97\% нь дээд үнэ нь доод үнээс бага буюу байж боломгүй лаа байсан.
\end{itemize}

\textbf{Шийдвэр:} FLF-LSTM загварыг суурь болгон ашиглахгүй. Энэ судалгаанд загвар бүрийг заавал энгийн жишиг
(baseline) аргатай харьцуулж, test өгөгдлийг зөвхөн эцсийн үнэлгээнд нэг удаа ашиглана.

\section{Бусад холбогдох судалгаа (10-р сар, 3-р долоо хоног)}
\begin{itemize}
	\item Sezer нарын тойм \cite{sezer2020financial}: ханш, хувьцааны үнийг гүн сургалтаар таамагласан 140 гаруй
	судалгааг шалгаад ихэнх нь дараах хоёр алдаа гаргадгийг тогтоосон.
	\begin{itemize}
		\item \textbf{Жишигтэй харьцуулаагүй} (naive benchmark): загвараа хамгийн энгийн аргатай, жишээ нь
		"маргаашийн ханш өнөөдрийнхтэй адил" гэсэн санамсаргүй алхалтын (random walk) таамаглалтай
		харьцуулдаггүй.
		\item \textbf{Арилжааны зардлыг тооцоогүй} (transaction costs): арилжаа бүрт төлөх spread, шимтгэлийг
		тооцдоггүй тул цаасан дээр ашигтай мэт харагддаг.
	\end{itemize}
	\emph{Дүгнэлт:} загвар бүрийг заавал жишиг аргатай харьцуулж, арилжааны зардлыг тооцно.
	\item Andersen нар \cite{andersen2003micro}: мэдээний бодит утга хүлээлтээс зөрсөн хэмжээ (surprise) ханшийг
	хэдхэн минутын дотор хөдөлгөдөг. \emph{Дүгнэлт:} мэдээний цаг минутын нарийвчлалтай зөв байх ёстой, мэдээний
	нөлөөг тусад нь шинжилнэ.
	\item LightGBM \cite{ke2017lightgbm}: хүснэгт хэлбэрийн өгөгдөлд хурдан, тогтвортой, шинжийн нөлөөг тайлбарлах
	боломжтой. \emph{Шийдвэр:} 12-р сард LightGBM-ийг гол загвар болгож, LSTM-тэй харьцуулна.
\end{itemize}
